\documentclass[10pt,a4paper]{amsart}
\usepackage[margin=19mm]{geometry}
\usepackage{amsmath,amssymb,mathtools}
\usepackage[scr=rsfs,cal=euler]{mathalfa}
\usepackage{xfrac}
\usepackage[unicode,hidelinks,bookmarks=false]{hyperref}
\newcommand{\ie}{i.e.\ }
\newcommand{\cf}{cf.\ }
\newcommand{\resp}{respectively\ }
\newcommand{\Subsection}{\subsection}
\providecommand{\nobreakdash}{\nobreak\hbox{-}}
\setlength{\emergencystretch}{2em}
\multlinegap0pt
\usepackage{bm}
\usepackage{bbm}
\usepackage{dsfont}
\usepackage{xspace}
\usepackage{cancel}
\usepackage{mathtools}
\usepackage{amsthm}
\makeatletter
\@ifundefined{defn}{\newtheorem{defn}{Defn}}{}
\@ifundefined{theorem}{\newtheorem{theorem}{Theorem}}{}
\@ifundefined{claim}{\newtheorem{claim}[theorem]{Claim}}{}
\makeatother
\makeatletter
\newcommand{\F}{\mathbb{F}}
\newcommand{\Span}{\mathop{\mathrm{Span}}}
\expandafter\ifx\csname claimproof\endcsname\relax
\newenvironment{claimproof}[1]{\par\noindent{\em Proof of Claim.}~#1}
\fi
\newcommand{\sympl}[1]{\langle #1 \rangle}
\makeatother
\title[Aspects of the commuting graph]{Aspects of the commuting graph}
\author{V. Arvind, Xuanlong Ma, Peter J. Cameron, Natalia V. Maslova}
\date{Source: arXiv:2305.07301v5 (CC BY 4.0)}
\begin{document}
\maketitle

\noindent\emph{Source and licence.} Aspects of the commuting graph. Version arXiv:2305.07301v5 (CC BY 4.0), distributed under the Creative Commons Attribution 4.0 International licence, as recorded in the arXiv licence metadata for this exact version and shown on the abstract page https://arxiv.org/abs/2305.07301v5. Full rights notice: licenses/arxiv-2305.07301-CC-BY-4.0.txt.\par\smallskip

\noindent\textit{Editorial scope.} Counted unit: the Theorem whose source label is \texttt{ortho-recog-thm} in the article (source0.tex lines 1419--1426, source label \texttt{ortho-recog-thm}), delivered together with the article's own complete original proof of it (source0.tex lines 1428--1747, one \texttt{proof} environment of 320 lines). No mathematics is written, repaired or re-derived: statement and proof are the article's own text line for line. Required context retained uncounted: sympl-basis-def (lines 1177--1187). Every definition, display and label of those lines is retained unchanged. Omitted from this package and not redistributed: 1--1176, 1188--1418, 1427--1427, 1748--4015. Nothing inside a retained excerpt is omitted or altered: every retained body is delivered exactly as the article's lines are written, so no omission receipt of any kind accompanies this package. Citations: the retained excerpts carry no citation command at all, so no bibliography entry of the article is transcribed and no reference is redistributed. Reference adaptation: none was needed and none was made; every reference inside the delivered unit resolves inside it, so no unresolved reference or citation is produced. Printed slips kept literal: no printed word, symbol, hypothesis or number of the counted statement or of its proof was altered. Layout and class: the article's own class is \texttt{elsarticle} with packages including amsmath, amsfonts, amssymb, url, algorithm, algpseudocode, xcolor; the wrapper is the lane's pinned \texttt{amsart} editorial wrapper at 10pt on A4, which restates the theorem-like environments the retained text needs and adds the article's own macro definitions that the retained text needs (\texttt{\textbackslash F}, \texttt{\textbackslash Span}, \texttt{\textbackslash claimproof}, \texttt{\textbackslash sympl}), with extra wrapper packages (bm, bbm, dsfont, xspace, cancel, mathtools). The delivered numbering is the unit's own. Assets: no retained span contains a figure environment, a graphics-inclusion command, a PDF-page inclusion, a table or a TikZ picture; no image file is copied into this package, the package declares no asset dependency and no image file is redistributed with it. Licence: version arXiv:2305.07301v5 of this article is distributed under the Creative Commons Attribution 4.0 International licence, as recorded in the arXiv licence metadata for this exact version and shown on the abstract page https://arxiv.org/abs/2305.07301v5. A full rights notice is delivered at licenses/arxiv-2305.07301-CC-BY-4.0.txt. Characters: every retained excerpt is pure ASCII, so the delivered engine, XeLaTeX with its default Latin Modern text font, reports no missing character.
\begin{defn}\label{sympl-basis-def}
  A \emph{symplectic basis} for the symplectic space $V(2r,p)$ is a
  basis of $2r$ vectors $\bar{x}_i, \bar{y}_i, 1\le i\le r$ such
  that:
\begin{itemize}
\item $\sympl{\bar{x}_i,\bar{y}_i}=1$ for all $i$.
\item $\sympl{\bar{x}_i,\bar{y}_j}=0$ for all $i\ne j$.
\item $\sympl{\bar{x}_i,\bar{x}_j}=0$ for all $i,j$.
\item  $\sympl{\bar{y}_i,\bar{y}_j}=0$ for all $i,j$.
\end{itemize}
\end{defn}

\begin{theorem}\label{ortho-recog-thm}
  Given a simple undirected graph $X=(V,E)$ on a vertex set $V$ of
  size $p^{2r}$, for an odd prime $p$, in time polynomial in the size
  of $X$ we can recognize if $X$ is the associated symplectic graph of
  an extraspecial group of order $p^{2r+1}$. Furthermore, the
  algorithm will also compute a bijection from $V$ to $V(2r,p)$ and
  determine an associated symplectic form $\sympl{,}$.
\end{theorem}

\begin{proof}
  The somewhat detailed proof is structured as a series of claims
  followed by a description of the algorithm and its analysis.

  Let $X=(V,E)$ be a graph with $p^{2r}$ vertices. Our algorithm will
  be recursive, with correctness proof based on an inductive
  argument. We know that the symplectic space $V(2r,p)$ has a
  symplectic basis of $2r$ vectors of the form
  $\bar{x}_i, \bar{y}_i, 1\le i\le r$ as defined in
  Definition~\ref{sympl-basis-def}.

%\begin{itemize}
%\item $\sympl{\bar{x}_i,\bar{y}_i}=1$ for all $i$.
%\item $\sympl{\bar{x}_i,\bar{y}_j}=0$ for all $i\ne j$.
%\item $\sympl{\bar{x}_i,\bar{x}_j}=0$ for all $i,j$.
%\item  $\sympl{\bar{y}_i,\bar{y}_j}=0$ for all $i,j$.
%\end{itemize}

We first note that given a symplectic form $\sympl{,}$ on $V(2r,p)$,
we can construct a symplectic basis for $V(2r,p)$ with the following
greedy procedure.
\begin{itemize}
\item Pick $\bar{x}_1\ne 0$ in $V(2r,p)$ arbitrarily.
\item Since $\sympl{,}$ is non-degenerate, we can find a vector
  $\bar{u}\in V(2r,p)$ such that $\sympl{\bar{x}_1,\bar{u}}=\alpha\ne
  0$.  Define $\bar{y}_1 =\alpha^{-1}\bar{u}$ such that
  $\sympl{\bar{x}_1,\bar{y}_1}=1$.
\item Then notice that the subspace $\{\bar{x}_1,\bar{y}_1\}^\perp$ is
  a $2r-2$ dimensional symplectic space $V'$ (w.r.t. the same
  symplectic form). We can continue with the basis construction,
  recursively applying the above two steps to $V'$.
\end{itemize}

Now, suppose $X=(V,E)$ is the graph given as input on $p^{2r}$
vertices. If $X$ is indeed the symplectic graph of $V(2r,p)$ for some
symplectic form $\sympl{,}$, we will describe a polynomial-time
recursive procedure that recovers a bijection from $V$ to $V(2r,p)$
along with a symplectic form $\sympl{,}$ as witness for $X$ being the
associated symplectic graph. We will then argue that if $X$ is not a
symplectic graph then this procedure detects that.

Essentially, the problem is to efficiently simulate the above two-step
basis construction procedure given only the graph $X=(V,E)$ as input.
As the symplectic space is non-degenerate, the zero vector $0$ is the
only vertex in $X$ adjacent to all others and is easily identified.
Let $x_1$ be any other vertex in $V$ and we choose a non-zero vector
$\bar{x}_1\in V(2r,p)$ as its image. Next, we will choose $y_1$ as any
vertex not adjacent to $x_1$ in the graph $X$ and choose a non-zero
vector $\bar{y}_1\in V(2r,p)$ as its image. Let

\[
V'=\{v\in V\mid (v,x_1)\in E \text{ and } (v,y_1)\in E\}.
\]
That is, $V'$ is the common neighborhood of $x_1$ and $y_1$.  We can
easily identify $V'$ in $X$. Let $X'$ be the subgraph of $X$ induced
by the vertex subset $V'$. We have the following simple observation.
\begin{claim}\label{symp-cl1}
 If $X$ is the symplectic graph of a $2r$-dimensional symplectic space
 over $\F_p$ then the graph $X'$ is the symplectic graph of a
 symplectic space of dimension $2r-2$ over $\F_p$.
\end{claim}

The algorithm will now recursively check that $X'$ is indeed the
symplectic graph of a $2r-2$-dimensional symplectic space $V(2r-2,p)$
over $\F_p$. If so, it will compute a labeling of the vertices of $V'$
by linear combinations $\sum_{i=2}^r(\alpha_i \bar{x}_i + \beta_i
\bar{y}_i), \alpha_i,\beta_i\in \F_p$ that is consistent with the
orthogonality relation of a symplectic form $\sympl{,}$. Here the
basis vectors computed by the algorithm are $\bar{x}_i,\bar{y}_i$ for
$2\le i\le r$. At the end of the recursive call, the algorithm computes a
bijection from $V'$ to $V(2r-2,p)$ that labels each vertex in $V'$ by
a distinct linear combination of the basis vectors.

Now, the remaining task for the algorithm is to find a labeling of the
vertices in $V\setminus V'$ consistent with a symplectic form on
$V(2r,p)$, assuming $X$ is a symplectic graph.

For each vertex $v$ we will denote its vector label by $\bar{v}$.  For
a labeled subset $S$ of vertices of $V$ we will use $\bar{S}\subseteq
V(2r,p)$ to denote the set of vectors labeling vertices in $S$.

\begin{claim}\label{symp-cl2}
  A vertex $v\in V\setminus V'$ can be labeled by a non-zero vector
  $\alpha \bar{x}_1 + \beta \bar{y}_1$ in $V(2r,p)$,
  $\alpha,\beta\in\F_p$, if and only if $(v,u)\in E$ for all
  $u\in V'$.
\end{claim}

\claimproof{Consider the symplectic graph of $V(2r,p)$. Let
  $V(2r-2,p)$ denote the subspace spanned by $\bar{x}_i,\bar{y}_i, 2\le i\le n$.
  Clearly, every vector of the form $\alpha \bar{x}_1 + \beta \bar{y}_1,
  \alpha,\beta\in\F_p$ is orthogonal to each vector in
  $V(2r-2,p)$. Conversely, consider a vector $\alpha \bar{x}_1 + \beta \bar{y}_1
  +v\in V(2r,p)$, where $v\in V(2r-2,p)$ is non-zero. Since $V(2r-2,p)$
  is non-degenerate, there is a $u\in V(2r-2,p)$ such that
  $\sympl{v,u}\ne 0$ which implies $\sympl{\alpha \bar{x}_1 +\beta \bar{y}_1 +
    v,u}=\sympl{v,u}\ne 0$.}

Thus, the algorithm can find a subset $V_{12}$ of precisely $p^2-1$
many such vertices in $V$ that are adjacent to all of $V'$, of which
we have already labeled two vertices as $\bar{x}_1$ and $\bar{y}_1$.
We analyze the subgraph $X_{12}$ of $X$ induced by $V_{12}$.

\begin{claim}\label{symp-cl3}
  The subgraph $X_{12}$ induced by $V_{12}$ consists of a disjoint
  union of $p+1$ many $(p-1)$-size cliques,
  $C_1,C_2,\ldots,C_{p+1}$. Furthermore, the labeling procedure must
  label the vertices of each $C_i$ from the set
  $\{\alpha \bar{u}\mid \alpha\in\F^*_p\}$ of $p-1$ vectors, for some
  non-zero vector $\bar{u}\in\Span(\bar{x}_1,\bar{y}_1)$.
\end{claim}

\claimproof{It is clear from the definition of $V_{12}$ that the
  vertices in it must be labeled by nonzero vectors in
  $\Span(\bar{x}_1,\bar{y}_1)$. The proof follows easily by observing
  that a non-zero vector $\alpha\bar{x}_1+\beta\bar{y}_1$ is
  orthogonal only to itself and its scalar multiples and to no other
  vectors in $\Span(\bar{x}_1,\bar{y}_1)$.}

By the above claim, we can consistently label each vertex $u\in
V_{12}$ by a vector $\bar{u}\in\Span(\bar{x}_1,\bar{y}_1)$ assuming
$X$ is a symplectic graph.

\begin{claim}\label{symp-cl4}
  If $X=(V,E)$ is the symplectic graph of the symplectic space
  $V(2r,p)$ and $V'$ corresponds to the symplectic subspace
  $V(2r-2,p)$ spanned by $\{\bar{x}_i,\bar{y}_i\mid 2\le i\le r\}$
  then for any vertex $u\in V_{12}$ the vertex subset
  \[
  V[u]=\{v\in V\setminus V'\mid (v,u)\in E\}
  \]
has to be labeled by the vectors
\[
\{\alpha \bar{u} +\bar{v}\mid \alpha\ne 0,
\bar{v}\in\Span(\cup_{i=2}^r\{\bar{x}_i,\bar{y}_i\})\}
\]
of $V(2r,p)$.
\end{claim}

\claimproof{Suppose $v\in V\setminus V'$ is labeled by the vector
  $\bar{a}+\bar{b}$, where $\bar{a}\in\Span(\bar{x}_1,\bar{y}_1)$ and
  $\bar{b}\in V(2r-2,p)$. As $v\in V[u]$ we have
  \[
  \sympl{\bar{a}+\bar{b},\bar{u}}=0.
\]
  As $\sympl{\bar{b},\bar{u}}=0$ it implies that
  $\sympl{\bar{a},\bar{u}}=0$ which means $\bar{a}=\alpha
  \bar{u}$ for some non-zero scalar $\alpha$. Conversely, any
  vertex labeled by $\alpha \bar{u} +\bar{b}$ for $\bar{b}\in
  V(2r-2,p)$ has to lie in $V[u]$ if the labeling is consistent with
  the symplectic graph $X$.}

It follows that $|V[u]|=(p-1)\cdot p^{2r-2}$. Furthermore, if vertices
$u$ and $v$ are in the same $(p-1)$-clique in the graph $X_{12}$ then
$V[u]=V[v]$. As
$(p+1)(p-1)\cdot p^{2r-2}=p^{2r}-p^{2r-2}=|V\setminus V'|$, it follows
that the different $V[u]$ are all pairwise disjoint and form a
partition of $V\setminus V'$. Consequently, we can identify the $p+1$
many equisized vertex subsets $V[u]$ which form a partition of
$V\setminus V'$.

\subsubsection*{Labelling vertices in $V[u]$}

By Claim~\ref{symp-cl4}, the vertices of $V[u]$ have to be labeled by
vectors from $\{\alpha\bar{u}+\bar{v}\mid v\in V', \alpha\ne 0\}$.  We
will examine the adjacencies in the subgraph induced by $V[u]$ and
also the adjacencies across the different $V[u]$ and $V[u']$ to
determine a consistent labeling. Let $X[u]$ denote the subgraph of $X$
induced by $V[u]$.

For distinct vertices $v_1, v_2\in V'$ notice that
\begin{equation}\label{eqn-V[u]}
  \sympl{\bar{v}_1,\bar{v}_2} =\sympl{\alpha \bar{u}+\bar{v}_1,\beta
    \bar{u}+\bar{v}_2}.
\end{equation}
Thus, in any consistent labeling of $V[u]$ suppose
$S_{u,v_i}\subset V[u]$ are the $p-1$ vertices labeled by
$\{\alpha\bar{u}+\bar{v}_i\mid \alpha \ne 0\}$ for $i\in\{1,2\}$. Then
Equation~\ref{eqn-V[u]} implies that in the subgraph $X[u]$ the
bipartite graph induced by the subsets $S_{u,v_1}$ and $S_{u,v_2}$
is either complete or empty.

Furthermore, $\sympl{\alpha\bar{u}+\bar{v},\beta\bar{u}+\bar{v}}=0$
for all non-zero $\alpha,\beta$. Hence each $S_{u,v}, v\in V'$ induces
a $(p-1)$-clique.

How do we identify each $S_{u,v}$? Suppose $v_1,v_2\in V'$ are
distinct vertices. In $V'$ (which is already labeled by $V(2r-2,p)$)
we can find a vertex $v_3$ such that $\sympl{v_1,v_3}\ne 0$ and
$\sympl{v_2,v_3}=0$. Hence each $S_{u,v}$ can be identified because
\begin{itemize}
\item All $p-1$ vertices in $S_{u,v}$ are adjacent to each other. They
  all have identical neighborhoods in $X[u]$ by the above argument.
\item Every vertex in $V[u]\setminus S_{u,v}$ has a different
  neighborhood in $X[u]$ than the vertices in $S_{u,v}$.
\end{itemize}

Putting it together, in summary, we have shown the following.

\begin{claim}\label{symp-cl5}
  In polynomial time we can uniquely identify all the subsets
  $S_{u,v}, v\in V'$, each of size $p-1$, forming a partition of
  $V[u]$, for each $u\in V_{12}$.
\end{claim}

Thus, we have identified the subsets $S_{u,v}\subset V[u]$ for $v\in
V'$. We now need to label the elements of $S_{u,v}$ by vectors
$\{\alpha\bar{u}+\bar{v}\mid \alpha\in \F^*_p\}$.

\begin{claim}\label{symp-cl6}
Assuming $X$ is a symplectic graph, for each $v\in V'$ and $u\in
V_{12}$, we can label the $p-1$ vertices in $S_{u,v}$ by the $p-1$
vectors $\alpha\bar{u}+\bar{v}, \alpha\in\F^*_p$ arbitrarily to get a
labeling for $V\setminus V'$ that is consistent with the labeling of
$V'$.
\end{claim}

\claimproof{We first observe that adjacencies within $V[u]$ and
  between $V[u]$ and $V'$ are consistent with such a labeling by the
  arguments preceding the claim.

  We now examine the adjacencies across two disjoint $V[u]$ and
  $V[u']$. Consider $S_{u,v}\subset V[u]$ and $S_{u',v'}\subset
  V[u']$. There is an edge between these subsets if and only if for
  some non-zero scalars $\alpha$ and $\beta$ we have
  $\sympl{\alpha\bar{u}+\bar{v},\beta\bar{u'}+\bar{v'}}=0$. Now,
  $\sympl{\bar{u},\bar{u'}}\ne 0$. Hence we must have
  \[
    \alpha\beta = \frac{-\sympl{\bar{v},\bar{v'}}}{\sympl{\bar{u},\bar{u'}}}.
  \]
  Since $\alpha$ and $\beta$ are non-zero scalars, it follows that
  $\sympl{\bar{v},\bar{v'}}\ne 0$. It follows that $v$ and $v'$ being
  non-adjacent in the graph induced by $V'$ is a necessary condition
  for $S_v$ and $S_{v'}$ to have any edges between them. Furthermore,
  the above equation implies that for the vertex, say $v_\alpha\in
  S_{u,v}$ labeled $\alpha\bar{u}+\bar{v}$ there is a unique $\beta\ne
  0$ such that $v_\alpha$ is adjacent to the vertex, say $v'_\beta$ in
  $S_{u',v'}$ labeled $\beta\bar{u'}+\bar{v'}$. Hence, either there
  are no edges between $S_v$ and $S_{v'}$ (if $v$ and $v'$ are
  adjacent) or there is a perfect matching between $S_v$ and $S_{v'}$(
  if $v$ and $v'$ are non-adjacent).

  Since these are all the possible edges in $V\setminus V'$ and it is
  clearly consistent with any labeling of $S_{u,v}$ by the $p-1$
  vectors $\alpha\bar{u}+\bar{v}, \alpha\ne 0$ for all $u\in V_{12}$
  and $v\in V'$, this completes the proof of the claim.}

Putting it together, clearly if $X$ is the symplectic graph of some
non-degenerate symplectic form on $V(2r,p)$ the above algorithm
verifies that by constructing a symplectic form consistent with the
symplectic graph. The algorithm is described as a recursive procedure
below. Essentially, given a graph $X=(V,E)$ with $|V|=p^{2r}$ vertices
as input it will try to compute a labeling of $V$ by the $p^{2r}$
vectors in
$V(2r,p)=\{\sum_{i=1}^r(\alpha_i\bar{x}_i+\beta_i\bar{y}_i)\mid
\alpha_i,\beta_i\in \F_p\}$, such that the edges of $X$ are consistent
with some symplectic form on $V(2r,p)$. If $X$ is a symplectic graph
it will succeed. Otherwise, the procedure will fail and detect that
$X$ is not symplectic.

\medbreak

\noindent\textbf{Procedure Label$(X,2r,p)$}
\begin{enumerate}
\item[] \emph{The procedure takes as input a simple graph $X=(V,E)$ with
$p^{2r}$ vertices for an odd prime $p$ and computes a
labeling consistent with some symplectic form on $V(2r,p)$ if
$X$ is a symplectic graph. Otherwise, it returns ``fail''.}

\item Identify the unique dominant vertex $v_0\in V$ and label it by
  $0$ (the zero vector).

\item Pick any pair of nonadjacent vertices $x_1, y_1\in V\setminus
  \{v_0\}$ and label them by the basis vectors $\bar{x}_1$ and
  $\bar{y}_1$ respectively.

\item Let $V'$ be the set of all vertices adjacent to both $x_1$ and
  $y_1$. Check $|V'|=p^{2r-2}$. Let $X'$ be the subgraph induced by
  $V'$. Recursively call Label$(X',2r-2,p)$ and obtain a labeling of
  $V'$ by vectors in $V(2r-2,p)$ defined by the basis
  $\bar{x}_i,\bar{y}_i, 2\le i\le p$ and consistent with a symplectic
  form on $V(2r-2,p)$ (Claim \ref{symp-cl1}).

\item Determine the $p^2-1$ vertices $V_{12}$ (Claim \ref{symp-cl2})
  consisting of vertices in $V\setminus V'$ that are adjacent to all
  vertices in $V'$. Determine the partition of $V_{12}$ into $p+1$
  many $(p-1)$-cliques in the induced subgraph $G_{12}$ (Claim
  \ref{symp-cl3}).

\item From the adjacencies between $V\setminus V'$ and $V'$ determine
  the partition of $V\setminus V'$ into $p+1$ subsets $V[u]$
  (Claim \ref{symp-cl4}).

\item From the subgraphs $G[u]$ induced by the respective $V[u]$,
  determine the partition of $V[u]$ into $p-1$-subsets $S_{u,v}, v\in
  V'$ (Claim \ref{symp-cl5}).

\item Label the vertices in $S_{u,v}, v\in V'$ and $u\in V_{12}$ by
  the vectors $\alpha\bar{u}+\bar{v}, \alpha\ne 0$
  (Claim~\ref{symp-cl6}).

\item If any step above does not satisfy the stated property then
  output \emph{fail}. Else output \emph{pass} along with the complete
  labeling of $V$ by elements of $V(2r,p)$ as the witness that $X$ is
  a symplectic graph.
\end{enumerate}

\noindent\textbf{Running time analysis}~~Let $T(n)$ denote the
running time for constructing a symplectic form $\sympl{,}$ for
$V(2r,p)$ consistent with $X$ as the symplectic graph, where
$n=p^{2r}$ is the number of vertices in $X$. The inductive
construction and the rest of the computation implies the recurrence
$T(n) = O(|V|^2) +T(n-1)$, which gives an overall cubic bound
$T(n)=O(|V|^3)=O(p^{6r+3})$ on the running time.

Putting it together, we have a polynomial-time algorithm that checks
if $X$ is the symplectic graph for the symplectic space by finding
out a labeling of vertices by the vectors along with a consistent
symplectic form. \qed
\end{proof}

\end{document}
